\section*{Adaptaciones a CRISP-DM}
\addcontentsline{toc}{section}{Adaptaciones a CRISP-DM}

En esta fase hicimos los siguientes ajustes a CRISP-DM.

\begin{itemize}
  \item Juntamos los cuatro entregables de la fase 2 (reporte de recolección inicial, de descripción, de exploración y de calidad de los datos) en este documento, en secciones separadas.
  \item Ordenamos el reporte alrededor del BCS, que es el único objetivo del proyecto.
  \item Los datos de ordeño se quedan como contexto, porque solo 12 de las 79 vacas con BCS tienen datos de ordeño. Se resumen dentro de este mismo documento.
  \item Analizamos los datos de BCS y las fotografías en un notebook de Python que se llama \nolinkurl{01_du_bcs_problem.ipynb} y está en la carpeta \nolinkurl{src/notebooks/02-data-understanding/} del repositorio.
  \item El reto académico pide evidenciar el uso de PySpark y AWS en el procesamiento de los datos. Como el volumen real de los datos es pequeño (las fotos originales pesan 2.18 GB con copias y 1.80 GB sin copias), documentamos y justificamos la decisión en el módulo de Big Data del repositorio, cuyo reporte está en la carpeta \nolinkurl{reports/bigdata/}.
  \item Los riesgos que encontramos en los datos se evalúan con probabilidad, impacto y magnitud, con la misma escala del plan de riesgos de la fase 1.
\end{itemize}
